
\chapter{Runtime и порядок исполнения}
\label{chap:runtime}

\code{Runtime} отвечает за исполнение программы CIAO.
При создании он компилирует декларативную модель,
проверяет её корректность и подготавливает независимые
контексты автоматных объектов.

После этого Runtime может принимать команды, выполнять
переходы и сохранять историю исполнения.

\section{Обработка команд}

Поступающие команды помещаются в очередь FIFO:
первая добавленная команда обрабатывается первой.

На каждом шаге Runtime извлекает одно сообщение,
выбирает подходящий переход и выполняет его.

\begin{figure}[ht]
\centering
\begin{tikzpicture}[
  box/.style={
    draw=ciaoRule,
    fill=ciaoCodeBg,
    rounded corners=1mm,
    minimum height=19mm,
    text width=30mm,
    align=center,
    font=\small\sffamily
  },
  flow/.style={
    -{Stealth[length=2.5mm]},
    draw=ciaoBlue,
    thick
  },
  node distance=6mm
]
  \node[box] (take) {
    \textbf{Извлечь}\\
    сообщение из FIFO
  };

  \node[box,right=of take] (choose) {
    \textbf{Выбрать переход}\\
    \footnotesize
    состояние, операция, Guard
  };

  \node[box,right=of choose] (actions) {
    \textbf{Выполнить}\\
    Actions по порядку
  };

  \node[box,right=of actions] (commit) {
    \textbf{Зафиксировать}\\
    состояние и исходящие команды
  };

  \draw[flow] (take) -- (choose);
  \draw[flow] (choose) -- (actions);
  \draw[flow] (actions) -- (commit);
\end{tikzpicture}
\caption{Один шаг исполнения Runtime}
\label{fig:runtime-step}
\end{figure}

При выборе перехода Runtime учитывает текущее
состояние автомата, полученную команду и, если они
объявлены, условия Guard.

Команды, отправленные через \code{Emit} или \code{Emits},
накапливаются во время выполнения Actions и попадают
в очередь только после успешного завершения перехода.

Таким образом, отправка команды не означает её
немедленного исполнения. Каждое сообщение обрабатывается
отдельным шагом Runtime.

\section{Пошаговое и автоматическое исполнение}

Для обработки очереди Runtime предоставляет два метода:

\begin{itemize}
  \item \code{step()} обрабатывает не более одного сообщения.
  \item \code{run()} выполняет шаги последовательно,
    пока очередь не опустеет или не будет достигнут
    установленный лимит.
\end{itemize}

Вернёмся к примеру Lamp. Отправим команду включения
и выполним один шаг:

\begin{lstlisting}[language=Python,numbers=none]
program, lamp = build_program()
runtime = Runtime(program)

runtime.send(lamp, Lamp.switch_on)
result = runtime.step()
\end{lstlisting}

Метод \code{send} добавляет команду в очередь,
а \code{step} обрабатывает её и возвращает
\code{StepResult} --- результат одного шага.

Он содержит сведения о получателе, команде, исходном
и новом состояниях, выбранном переходе и результате
исполнения.

Проверим выполненный переход:

\begin{lstlisting}[language=Python,numbers=none]
assert result.transition == "turn_on"
assert runtime.snapshot(lamp).state is Lamp.on
\end{lstlisting}

Для обработки нескольких сообщений удобнее
использовать \code{run()}:

\begin{lstlisting}[language=Python,numbers=none]
runtime.send(lamp, Lamp.switch_off)
runtime.send(lamp, Lamp.switch_on)

runtime.run()
\end{lstlisting}

Лампа последовательно выключится и включится.
Если во время переходов автоматы отправят новые
команды, Runtime также добавит их в очередь
и продолжит исполнение.

Чтобы ограничить число шагов, можно передать
параметр \code{max_steps}:

\begin{lstlisting}[language=Python,numbers=none]
runtime.run(max_steps=100)
\end{lstlisting}

По умолчанию \code{run} уже имеет конечный лимит.
Он защищает программу от бесконечного исполнения,
например при циклической отправке командами самих себя.

\section{Ошибки исполнения}

В примере Lamp мы уже пытались выключить лампу,
находящуюся в состоянии \code{off}.

Поскольку подходящий переход отсутствует,
Runtime вызывает исключение \code{NoTransitionError}.

По умолчанию Runtime работает в строгом режиме.
Если указать \code{strict=False}, сообщения без
подходящего перехода будут пропускаться, а их
результат получит статус \code{ignored}.

Нестрогий режим может быть полезен, когда программа
допускает получение команд, не предусмотренных
для текущего состояния автомата.

\begin{importantbox}
Нестрогий режим не подавляет ошибки пользовательских
Guard и Action.

Если во время Action возникает исключение, Runtime
не устанавливает целевое состояние и не публикует
команды, подготовленные этим переходом.

Однако уже выполненные изменения контекста
автоматически не отменяются.
\end{importantbox}

\section{Состояние и трасса исполнения}

Текущее состояние автоматного объекта и его контекст
можно получить через \code{snapshot}:

\begin{lstlisting}[language=Python,numbers=none]
snapshot = runtime.snapshot(lamp)

print(snapshot.state)
print(snapshot.context.switch_on_count)
\end{lstlisting}

Метод возвращает снимок данных, а не изменяемую
ссылку на внутренний контекст Runtime.

Для анализа последовательности переходов
Runtime также сохраняет трассу исполнения.

Каждая запись \code{TraceEvent} содержит информацию
об обработанном сообщении: получателя, операцию,
исходное и целевое состояния, выбранный переход
и результат выполнения.

В отличие от \code{StepResult}, который описывает
отдельный шаг, трасса позволяет проследить историю
исполнения программы.

Это особенно полезно при отладке систем из нескольких
автоматов, когда результат зависит от порядка
поступления и обработки команд.
